\section{Buffered circuits, mixing, and rare pins}\label{sec:cuts}

This section converts two loop estimates into the conditional comparisons
needed throughout the proof.  The progression is from a circuit under an
extreme one-color boundary condition, to circuits in fixed geometric
buffers, to a plane law and multiplicative estimates for rare pins.  At
each step the selected cut leaves its strict interior unexamined.

We first fix the geometry of the loop estimates imported from the literature.
Let $\mathbb H=\Tri^*$ be the honeycomb lattice.  Its hexagonal faces are
identified with the height vertices of $\Tri$; face distance means graph
distance in $\Tri$.  For a center face $x$, let $\Gamma_n(x)$ be the outer
honeycomb contour of the union of closed faces at distance at most $n$
from $x$.  Let $\Lambda_n(x)$ be the honeycomb domain strictly inside
that contour, excluding its vertices and edges, and write $\Lambda_n$
when $x=0$.  These are face balls in the honeycomb lattice, not balls of
elementary triangular faces.

The \emph{graph inradius} of a simple contour surrounding $x$ is the largest
integer $n$ for which $\Lambda_n(x)$ and its bounding contour lie strictly
inside it.  This differs by a universal additive constant from the minimum
graph distance to a face incident to the contour: the smaller connected
face balls cannot cross the contour, and the next incident face layer
reaches it.  Edge, vertex, and integer-rounding conventions change only
this bounded offset.

Glazman--Manolescu \cite[Theorem~1.2(iv), equations~(1.1)--(1.2)]{GM2021} prove the following facts
for the loop $O(2)$ model with edge weight one and empty exterior condition:
\begin{enumerate}[label=(\roman*)]
\item A domain containing $\Lambda_{2n}$ has probability at least $c_0>0$
of a loop contained in $\Lambda_{2n}$ and surrounding $\Lambda_n$.
\item There are $0<\rho<1$ and $c_0>0$ such that, if the graph inradius
about the center face is $n$, the probability of two loops surrounding
$\Lambda_{\rho n}$ is at most $1-c_0$.
\end{enumerate}
All radii in these estimates use the face-graph convention just defined.
The estimates are uniform for sufficiently large integer radii.  The
source's contour-distance convention differs by a bounded additive offset;
choose $\rho$ slightly larger than its upper-bound constant, still below
one, to absorb that offset and target-ball rounding for all large radii.
No comparison with Euclidean distance is used in this adjustment.  The
domains are interiors of simple honeycomb contours: loops are strictly inside
the contour and do not use its vertices.  When we use the second estimate
in a more general height domain, we first expose an enclosing level loop;
its interior has exactly this empty-exterior law by \cref{lem:gray}.

\begin{lemma}[A circuit under the opposite cut condition]\label{lem:worst-cut}
There are fixed $K<\infty$ and $c>0$ with the following property.
Inside a regular dual cut of radius $Kn$, fix $B=-t$ on the inner boundary
layer and leave $W$ free.  With probability at least $c$, there is a
double-$B$ circuit of sign $t$ between radii $n$ and $2n$.  The constants
are independent of $n$, for all sufficiently large $n$.
\end{lemma}

\begin{proof}
It is enough to take $t=+1$.  In a larger auxiliary flat height domain,
choose the constant boundary residue uniformly from the four Gray
residues.  Require enclosing loops in three separated annuli, with the
inner annulus $\Lambda_{2n}\setminus\Lambda_n$, and with the outer
annulus still strictly inside the prospective deterministic cut.  For
example, annuli at radii $n$, $4n$, and $16n$ fit inside a cut of a fixed
sufficiently large multiple of $n$.

To bound the probability, find an outer qualifying loop from the exterior
and then repeat inside it.  Each remaining target annulus is contained
with its required buffer in the free interior left by the preceding loop.
The lower loop estimate and \cref{lem:gray} give a probability at least
$c_0^3$ for the three-loop event.

Conditional on all loop geometry, list the enclosing loops in order.
Their jump signs are independent fair coins.  At any height residue, the
two possible jumps flip different spin colors.  Thus the sequence of
flipped colors is also a sequence of independent fair choices, conditional
on the initial height parity.  The initial $B$ sign is fair independently
of that parity.  There is consequently a fixed positive conditional
probability that an outer selected loop flips $W$ with unchanged $B=-1$
and that an inner selected loop flips $W$ with unchanged $B=+1$.
An intervening selected loop supplies a free color choice to set the
parity of the intervening $B$ flips.  Additional intervening loops do not
reduce this bound: condition on their color choices first.  A loop that
flips $W$ is itself a double-$B$ circuit of the unchanged $B$ sign.

Let $A$ be the desired inner plus-circuit event, and search from the
exterior for a minus double circuit $C$ in the outer annulus.  The
construction just given yields
$\PP(C\hbox{ is found},A)\geq c_1>0$.  The stopping cut law says that,
conditional on this search, the interior is the free-$W$, minus-$B$ cut
law inside $C$.  Write its probability of $A$ as $q_C$.

In the deterministic farther minus-cut law, condition additionally on
$B=-1$ on both layers of this fixed realized $C$.  Since $A$ is increasing,
\cref{lem:spin-fkg} says that this conditioning can only decrease its
probability.  By \cref{lem:double-cut}, the resulting inner probability is
$q_C$.  If $q$ is the desired probability in the deterministic cut law,
then $q\geq q_C$ for every such $C$.  Integrating over the auxiliary
stopping search gives $q\geq c_1$, as required.
\end{proof}

For clarity, a \emph{buffer} below means an open neighborhood in which
the local annuli, together with their larger worst-condition cuts, fit
with positive clearance.  Thus an obstacle can lie in the hole of a
collar without lying in the buffers used to construct the collar circuit.
We distinguish these local buffers from the entire disk bounded by the
eventual circuit.

\begin{lemma}[Circuits in buffers]\label{lem:buffer-circuits}
Consider a fixed positive-width polygonal collar, or a fixed positive-width
tube following a polygonal arc.  Suppose it has a buffer disjoint from
$E_0$ and from all prescribed $B$ spins.  At sufficiently fine mesh,
conditional on arbitrary $B$ values away from the buffer, there is a
probability at least $c>0$ of a double-$B$ circuit around the collar's hole,
or a connected double-$B$ path following the tube, of either specified sign.
For a sign-$t$ event, arbitrary additional $B=t$ pins in the buffer are
allowed.  The constant depends on the fixed geometry, not on the mesh,
the number of pins elsewhere, or their values.

Finitely many separated such requirements have a uniformly positive joint
probability, even if different buffers request different signs.  For a
logarithmic stack of separated similar annular buffers between radii $r$
and $L$, the probability that no circuit of a prescribed sign occurs is
at most
\begin{equation}\label{eq:stack-failure}
 C(r/L)^\alpha
\end{equation}
for some $\alpha>0$.  In lattice units this holds uniformly for $L\gg r\geq1$,
after adjustment of $C$.
\end{lemma}

\begin{proof}
First consider a small regular annulus whose larger cut lies in the
buffer.  Delete any favorable pins in that cut.  This can only lower the
probability of the desired event in the order toward $t$.  Fix the
opposite sign on both sides of the larger deterministic cut.  FKG lowers
the desired probability again, and the exact cut law identifies the
remaining problem with \cref{lem:worst-cut}.  This gives a uniform local
lower bound under every allowed exterior conditioning.

We explain the geometric gluing, including why a thin collar requires
no stronger imported loop estimate.  Cover the center curve of the tube
or collar by finitely many small regular annuli.  Choose their common
inner radius much smaller than the tube width.  Consecutive centers may
be separated by, for example, $3/2$ times that radius in the graph-distance
norm.  Their inner balls intersect, and each inner ball protrudes beyond
the other's outer ball.  Two surrounding circuits in these annuli must
therefore intersect: disjoint Jordan curves would have disjoint interiors
or nested interiors, and the two geometric conditions exclude these
possibilities.  Their union is connected and follows the prescribed
tube.  Around a closed center curve, its outer boundary contains a simple
circuit enclosing the hole; the annuli are small enough that their union
does not fill that hole.  At a common dual vertex, double paths of the
same sign have all three adjacent spin sites of that sign, so the gluing
does not lose the double-path property.  Bounded lattice rounding is
absorbed by the positive geometric clearances.

All the small annulus events for one sign are increasing in the same
order.  FKG multiplies their local lower bounds, giving a positive
constant depending only on the finite cover.  For separated buffers,
condition successively on the configurations outside the next buffer;
the preceding lower bound applies under each such conditioning.  This
also permits different signs in different buffers.

Finally choose a fixed scale ratio large enough that the enlarged buffers
of successive annular slots are disjoint.  There are at least
$c_2\log(L/r)-C_2$ slots.  At each trial, conditional on the previously
examined slots, success has probability at least $c_3>0$.  Failure in all
slots has probability at most $(1-c_3)^{c_2\log(L/r)-C_2}$, which is
\eqref{eq:stack-failure}.  The finitely many bounded lattice scales can be
absorbed in its prefactor.
\end{proof}

We record the coupling mechanism behind our use of these circuits.  Take
two pair laws whose domains contain a large ball, with no constraints
there.  On a slightly smaller deterministic dual cut, compare each $B$
marginal with the law having minus fixed on both boundary layers.
Exact cutting makes this lower interior law common to both comparisons.

Run the exterior exploration of \cref{lem:double-cut} in the lower law,
querying the same sites in the other two laws.  At every stage retain,
for each marginal, its conditional law given its own queried values.
If these values are ordered, conditional FKG orders the next site's plus
probabilities: fixing the outer minus layer and then lowering the queried
values can only decrease that probability.  A fresh uniform random
variable therefore samples all three next spins in their required order.
The next query depends only on the revealed record, so this recursive
conditional sampling preserves each marginal.

Stop when the lower configuration has a plus double circuit.  Its cut
sites are plus in all three configurations.  The exploration has queried
only their exterior and these cut sites, so their remaining conditional
interior laws coincide.  Fill that interior identically.  Sample $W$ only
after this $B$ exploration; its components are free across the cut and can
also be matched.  The construction allows arbitrary constraints outside
the large ball, including $E_0$ equalities there.

The same construction compares a law with one obtained by adding extreme
$B$ pins in a small region.  Explore outward from that region in the more
adverse configuration.  A double circuit of the added sign matches the
configurations on the far side.  This version will be used when filling
small gaps in a pin curve.

\begin{proposition}[Plane law and quantitative local comparison]
\label{prop:plane-mixing}
Free coherent pairs in expanding regular hexagons converge on cylinder
events to a law $\mu_\infty$.  This law is invariant under lattice
translations, lattice symmetries, global sign changes, and interchange
of the colors.  If the domain of a finite pair law $\mu^{E_0,\eta}$
contains a ball of radius $L$ and has no constraints in that ball, then
its restriction to a concentric
ball of radius $r$ differs from that of $\mu_\infty$ in total variation
by at most $C(r/L)^\alpha$, for $L\gg r\geq1$.

The associated plane height increments exist and have the same spatial
symmetries.  Every finite feasible prescription of $B$ pins together with
$W=+1$ on finitely many connected paths has positive plane probability.
\end{proposition}

\begin{proof}
Apply the preceding common-lower-law coupling and search for a plus circuit
in a stack between the inner ball and the deterministic outer cut.  Its
failure probability is bounded by \eqref{eq:stack-failure}.  On success
the full pair configurations agree in the inner ball.  This bounds the
total variation distance between any two sufficiently large free-volume
laws, and proves the Cauchy property on each cylinder.  It gives the same
bound with one constrained law and one expanding free volume; then let
the latter volume tend to the plane.  Constants do not depend on the
constraints outside the empty ball.

Translated or reflected exhausting volumes give the same limit by this
comparison, proving the spatial symmetries.  The finite laws have the
stated spin symmetries.  Formula \eqref{eq:height-increment} determines
increments on every fixed path, and its vanishing curl persists in the
limit.  It therefore defines the plane height modulo one additive
integer.

To check positivity of a finite pin event, enclose all its sites and paths
by a double-$B$ cut.  Such a cut has positive probability by
\cref{lem:buffer-circuits}.  In its finite interior, take $W$ everywhere
plus and choose $B$ with the prescribed values and the cut value.  This
is a coherent extension and has positive conditional counting probability.
The exact cut law proves the assertion.  Uniform finite-volume estimates
pass to this positive-probability plane conditioning.
\end{proof}

When height increments rather than the two spin labels are considered,
one may equivalently anchor a free height at zero.  The global Gray
rotation $(B,W)\mapsto(W,-B)$ preserves \eqref{eq:height-increment} and
cycles the four residues at the anchor.  Each anchored height consequently
has the same four representatives in the free pair law.  This observation
will let us apply counting inequalities directly to anchored height arrays.

\subsection{A comparison for separated rare pin events}

A family of \emph{separated obstacles} consists of finitely many supports,
each contained in a bounded simply connected polygonal neighborhood $U_i$,
with the closures $\overline U_i$ pairwise disjoint.  Each support has a
surrounding polygonal collar inside $U_i$, with positive clearance for all
its local buffers and larger comparison cuts.  Their filled interiors
also lie in $U_i$.  Thus an isolating cut encloses only its own obstacle;
disjoint annular collars alone would not impose this condition.
After rescaling, these neighborhoods, collars, and clearances are fixed.
For example, finitely many lattice paths converging uniformly to disjoint
compact simple arcs admit such neighborhoods: choose sufficiently thin
regular neighborhoods of the arcs, with endcaps, and approximate their
boundaries polygonally.  Their closures remain disjoint at fine mesh.  On obstacle $i$, an event $P_i$ prescribes arbitrary
$B$ values and requires $W=+1$ on at most $m_i$ connected subpaths.  There
is a fixed bound on $\sum_i m_i$, but no bound on the number of pinned
sites along a subpath.  We use the same notation $P_i$ for the event and
write $p_i=\mu_\infty(P_i)$ for its probability.

The following comparison treats one obstacle while retaining any pins
already imposed on it.  This retention allows its plus and minus pins to
be prescribed in successive batches.

\begin{lemma}[Relative comparison for one obstacle]
\label{lem:relative-pin-comparison}
Let $\mu$ and $\nu$ be finite pair laws with second-color equality
constraints and first-color pins.  Their domains contain the same filled
isolating neighborhood $U$, and their constraints inside $U$ agree:
these consist of retained $B$ pins $R$ and equality edges $E_0$ on one
obstacle.  Its surrounding collar and buffers in $U$ contain no
constraints.  Outside $U$ the two laws may have different constraints.
If $A$ prescribes one sign $t$ on further sites of that obstacle,
consistently with $R$, then
\[
 c\nu(A)\leq\mu(A)\leq c^{-1}\nu(A).
\]
Here $c>0$ depends only on the collar and its buffers, not on the
retained pins, equality edges, remote constraints, or the probability
of $A$.  The same assertion holds with the colors interchanged.
\end{lemma}

\begin{proof}
Fix sign $t$ on a deterministic double layer $\Gamma$ in the collar.
Let $q$ be the probability of $A$ in its interior law, retaining $R$
and the internal $E_0$ constraints.  Exact cutting makes $q$ the same
for $\mu$ and $\nu$.  FKG in the order toward $t$ gives, for either
law $\lambda\in\{\mu,\nu\}$,
\begin{equation}\label{eq:rare-pin-upper}
 \lambda(A)\leq\lambda(A\mid B=t\text{ on }\Gamma)=q.
\end{equation}
In a thinner collar strictly inside $\Gamma$, an exterior stopping search finds
a double circuit of sign $t$ with probability at least $c_1>0$, by
\cref{lem:buffer-circuits}.  Its constant is uniform over $R$ and all
remote constraints: the search buffers contain none of these pins or
$E_0$ edges.  Conditional on the search, the inner law is the corresponding
closer-cut law.  In the deterministic-$\Gamma$ law, conditioning on this closer
sign-$t$ cut can only increase $A$, while preserving $R$.  Its probability
of $A$ is therefore at least $q$.  Consequently
\begin{equation}\label{eq:rare-pin-sandwich}
 c_1q\leq\lambda(A)\leq q,
 \qquad \lambda\in\{\mu,\nu\}.
\end{equation}
Comparison through this common $q$ proves the lemma.  The retained pins
may be extremely rare; their probability never enters the inequalities.
The stopping exploration has not inspected the strict interior, which
is essential for the lower bound.
\end{proof}

\begin{proposition}[Products of separated pin probabilities]
\label{prop:pin-product}
For a separated family as above, there are constants $0<c\leq C<\infty$
such that
\begin{equation}\label{eq:pin-product}
 c\prod_i p_i\ \leq\ \mu_\infty\Bigl(\bigcap_iP_i\Bigr)
 \ \leq\ C\prod_i p_i.
\end{equation}
The constants depend on the collars, their buffers, the number of
obstacles, and $\sum_i m_i$.  They do not depend on the mesh, the number
or the values of the microscopic $B$ pins, or the probabilities $p_i$.
The same estimates hold in all sufficiently large finite free volumes.
\end{proposition}

\begin{proof}
First handle the equality requirements for $W$.  Temporarily
ignore all $B$ pins.  For a collection of $m$ connected paths, let $E$
denote constancy of $W$ along each path and let $L$ require all labels to
be plus.  Conditional on $B$, a consistent pattern of labels has the same
probability as the all-plus pattern, and an inconsistent pattern has
probability zero.  Summing over the at most $2^m$ patterns gives
\begin{equation}\label{eq:constancy-label-comparison}
 \PP(L)\leq\PP(E)\leq 2^m\PP(L).
\end{equation}
Apply \cref{lem:relative-pin-comparison} with the colors interchanged to the
successive all-plus requirements on the obstacles.  Each conditional
factor is comparable to the corresponding one-obstacle factor: the
conditioning concerns only other, separated obstacles.  The all-plus probability thus factors up to
constants.  Equation \eqref{eq:constancy-label-comparison}, for the joint
family and for its individual members, gives the same conclusion for
their constancy events.  Denote the resulting equality edges by $E_0$.

Now condition on all these equalities and impose the $B$ requirements
obstacle by obstacle.  On each obstacle first impose its plus pins, then
its minus pins.  At either stage compare its conditional factor with the
same stage in the one-obstacle equality problem using
\cref{lem:relative-pin-comparison}.  At the minus stage the plus pins already
imposed on that obstacle belong to $R$ and are retained in both cut laws.
All previous pins on other obstacles and all other equality paths lie
outside $U_i$, so the two laws have exactly the same retained constraints
inside $U_i$.  Multiplication compares the joint conditional $B$-pin
probability with the product of the one-obstacle conditional probabilities.

Finally restore all the required $W$-plus labels.  Conditional on $B$ and
$E_0$, their probability is between $2^{-\sum_i m_i}$ and one; the same
statement applies in each one-obstacle problem.  These bounded factors,
together with the preceding two comparisons, prove
\eqref{eq:pin-product}.  The entire proof takes place in finite volumes
with uniform constants.  The plane assertion follows from
\cref{prop:plane-mixing} and positivity of each finite pin event.
\end{proof}
